% =====================================================================
% MSc Preliminary Project Report
% MSc Computational Finance, King's College London
% Hüseyin Sarp Vulaş — Supervisor: Dr Bart de Keijzer
% Submission deadline: 30 April 2026
% =====================================================================
\documentclass[11pt,a4paper]{article}

% --- Encoding / fonts (handles ş, ü in author name) -------------------
\usepackage[utf8]{inputenc}
\usepackage[T1]{fontenc}
\usepackage{lmodern}

% --- Layout -----------------------------------------------------------
\usepackage[margin=2.0cm]{geometry}   % comfortably above the 1.5 cm minimum
\usepackage{microtype}
\setlength{\parskip}{4pt}
\setlength{\parindent}{14pt}

% --- Maths / tables / figures -----------------------------------------
\usepackage{amsmath,amssymb}
\usepackage{booktabs}
\usepackage{array}
\usepackage{graphicx}
\usepackage{caption}
\captionsetup{font=small,labelfont=bf}

% --- Citations --------------------------------------------------------
\usepackage[round,authoryear]{natbib}
\bibliographystyle{plainnat}

% --- Hyperlinks (load late) ------------------------------------------
\usepackage[hidelinks]{hyperref}

% --- Gantt chart ------------------------------------------------------
\usepackage{pgfgantt}

% --- Section formatting ----------------------------------------------
\usepackage{titlesec}
\titleformat{\section}{\large\bfseries}{\thesection}{0.6em}{}
\titleformat{\subsection}{\normalsize\bfseries}{\thesubsection}{0.6em}{}
\titlespacing{\section}{0pt}{10pt}{4pt}
\titlespacing{\subsection}{0pt}{6pt}{2pt}

% =====================================================================
\begin{document}

% --- Cover page (does not count toward 10-page limit) -----------------
\begin{titlepage}
\centering
\vspace*{2.5cm}
{\large King's College London \\ Department of Informatics \par}
\vspace{0.4cm}
{\large MSc Computational Finance, 2025/2026 \par}
\vfill
{\Large\bfseries Multi-Market Equity Crisis Early Warning \\[4pt]
with Regime-Aware Temporal Fusion Transformers \par}
\vspace{0.6cm}
{\itshape MSc Preliminary Project Report \par}
\vfill
{\large Hüseyin Sarp Vulaş \par}
\vspace{0.2cm}
{\normalsize Supervisor: Dr Bart de Keijzer \par}
\vfill
{\normalsize 30 April 2026 \par}
\vspace{1cm}
\end{titlepage}

% =====================================================================
% BODY STARTS HERE (10-page limit applies from this point)
% =====================================================================

\section*{Abstract}
\noindent
Single-market early warning systems (EWS) for equity crises suffer from a
structural data scarcity problem: a single index produces only a handful of
crisis episodes per decade, far too few for attention-based deep learning
architectures to learn meaningful temporal patterns. Recent empirical work has
documented that, in this regime, gradient-boosted trees match or beat
neural sequence models. We argue that this conclusion is an artefact of the
single-entity setting rather than a fundamental property of the problem, and
we propose to reformulate the task as a multi-market panel forecasting problem
spanning twelve major equity indices over twenty-four years.
On this richer dataset we develop the \emph{Regime-Aware Temporal Fusion
Transformer} (RA-TFT), which augments the Temporal Fusion Transformer of
\citet{lim2021tft} with a differentiable neural Hidden Markov Model and a
regime-conditioned attention bias. Preliminary experiments on a single
temporal split confirm the central hypothesis: on the panel dataset RA-TFT
attains a pooled ROC-AUC of $0.640$ against $0.622$ for XGBoost, reversing
the ranking observed on the single-market task. This report sets out the
problem, the methodology, the early empirical evidence, and the work plan
through to the 6 August 2026 final submission.

% =====================================================================
\section{Introduction}
\label{sec:intro}

\subsection{Domain and Problem}
Equity bear markets impose disproportionate costs on real-economy actors: in
the run-up to a crisis, capital is committed at peak valuations, while in the
aftermath credit conditions tighten precisely when firms and households most
need accommodation. A reliable \emph{early warning system} for crisis onset
is therefore of practical interest to risk managers, central banks, and
multi-asset allocators \citep{kaminsky1999leading,reinhartrogoff2009}. The
quantitative finance literature has accumulated a wide variety of EWS
proposals over the past quarter-century, but their performance on \emph{deep}
crisis events (those with drawdowns exceeding twenty per cent from peak)
remains underwhelming, particularly for sequence-aware architectures
that, on paper, ought to dominate this regime.

We focus on the binary classification of crisis onset, defined as the first
day on which a market enters a sustained drawdown episode of $\geq 20\%$ from
its trailing 252-day rolling peak. This labelling convention follows
\citet{dichtl2023ews} and is robust to short-lived corrections. The
prediction horizon is sixty-three trading days, corresponding to roughly one
calendar quarter, a horizon long enough to be actionable yet short enough
that the underlying signals (volatility regime, credit conditions,
cross-market contagion) remain informative.

\subsection{Aims and Objectives}
The primary aim of this project is to determine \emph{whether and under what
conditions} attention-based sequence architectures can outperform
gradient-boosted trees on equity crisis EWS, and, if so, to design and
empirically validate an architecture that exploits the structure of the
problem. Concretely we set the following objectives:

\begin{enumerate}
\item Reformulate the conventional single-market EWS task as a multi-market
panel forecasting problem with a sufficient number of crisis episodes to
support attention-based learning.
\item Develop the RA-TFT architecture, integrating a neural HMM regime
detection module and a regime-conditioned attention bias into the Temporal
Fusion Transformer of \citet{lim2021tft}.
\item Evaluate RA-TFT against four baselines (vanilla TFT, LSTM, XGBoost,
penalised logistic regression) using a crisis-anchored walk-forward
cross-validation protocol with bootstrap confidence intervals on PR-AUC.
\item Quantify the contribution of each architectural component
(regime module, regime attention, static market embedding) through a
controlled ablation.
\item Examine calibration quality (expected calibration error, Brier score)
and apply post-hoc methods (Platt scaling, isotonic regression) where
necessary, since EWS probabilities feed directly into downstream risk
management decisions.
\end{enumerate}

\subsection{Dataset}
The project uses daily price and volume data for twelve major equity indices
spanning North America, Europe, Asia-Pacific, and emerging markets, from
1 January 2000 through 31 December 2024 (24 years). The full panel
contains 72{,}801 market-days. Macro-financial covariates (VIX, the
VIX term structure, US high-yield and investment-grade credit spreads, the
TED spread, the dollar index, gold, and the MOVE index) are obtained from
FRED and Yahoo Finance and broadcast across all markets. The full instrument
list and the per-market and global feature definitions appear in
Section~\ref{sec:method}. All data are publicly available and ingestion is
already automated by the project pipeline.

\subsection{Methodological Approach}
RA-TFT inherits the variable selection networks, gated residual connections,
and interpretable multi-head attention of the original TFT, and supplements
them with two architectural elements designed for the crisis prediction
task. First, a neural HMM with three latent states (\emph{calm},
\emph{stress}, \emph{crisis}) and time-varying transition probabilities sits
on top of the LSTM encoder; its forward algorithm is differentiable end-to-end
and is initialised so that transitions are sticky (high diagonal probability)
which prevents the regime distribution from oscillating. Second, the
inferred regime distribution conditions a learned additive bias to the
attention logits in the temporal self-attention layer, allowing the model to
weight long-range historical context differently in calm versus stressed
periods.

\subsection{Contributions}
The contributions we expect to be in a position to claim by 6 August 2026 are
threefold. \emph{First}, an empirical demonstration that the choice between
attention-based and tree-based EWS models is not architectural in nature but
depends critically on the number of crisis episodes available: once the
single-market regime is left behind, the ranking inverts. \emph{Second}, a
novel RA-TFT architecture combining differentiable HMM regime inference
with regime-conditioned attention; to our knowledge no prior work has
integrated time-varying-transition neural HMMs into TFT for financial EWS.
\emph{Third}, a publicly reproducible pipeline (data ingestion through to
calibrated probabilistic forecasts) on a panel that other researchers can
extend.

% =====================================================================
\section{Background and Related Work}
\label{sec:background}

The relevant literature falls into four broadly overlapping strands: the
econometric origins of crisis early warning, the more recent application of
machine learning to crisis prediction, the development of attention-based
architectures for time series, and the long history of regime-switching
models in financial econometrics. We treat them synthetically rather than
paper-by-paper, since the contribution of this project lies at the
intersection.

\subsection{Crisis Early Warning Systems}
The modern EWS literature originates with the signals approach of
\citet{kaminsky1999leading}, who screened a large set of macroeconomic and
financial indicators for their ability to anticipate currency crises. This
approach evolved into a logistic-regression tradition through
\citet{frankelsaravelos2012}, who emphasised reserve adequacy and real
exchange rate misalignment, and into a deep-learning tradition through
\citet{beutel2019does}, \citet{holopainen2017toward}, and
\citet{chatzis2018forecasting}, who applied feed-forward and recurrent
networks to country-level crisis panels. A consistent finding across these
studies is that ensemble tree methods (random forests and gradient
boosting) match or beat neural alternatives, an observation crystallised
by \citet{bluwstein2023credit} on a long historical credit-cycle panel.
The implicit message is that crisis prediction is a low-data regime in which
the inductive bias of trees outperforms the flexibility of deep models.

A complementary methodological strand has refined the labelling problem
itself. Naïve drawdown regression mixes recovery from prior crises with
genuine crisis onset and produces noisy targets. \citet{dichtl2023ews} argue
for binary onset labelling with a fixed forward horizon and exclusion of
in-crisis observations from training, a convention we adopt here.

\subsection{Sequence Models in Financial Time Series}
Recurrent and attention-based architectures have a strong record in
forecasting tasks where data are abundant: high-frequency limit order book
prediction, volatility forecasting on long minute-bar histories, and
multi-horizon retail demand. Their record on crisis prediction, by contrast,
is weak, and we argue this is because the conventional single-country setup
exposes them to a fundamentally adverse data-to-parameter ratio. With
$\sim$4 crisis episodes in 25 years, an attention layer with $O(d^2)$
parameters has nothing to learn against. Two responses are possible:
shrink the model (trading representational capacity for sample efficiency,
which forfeits the reason for using attention in the first place), or
increase the effective number of crisis episodes through cross-entity
pooling. Our project pursues the second route.

\subsection{Temporal Fusion Transformer}
The Temporal Fusion Transformer of \citet{lim2021tft} is purpose-built for
panel forecasting. Its static covariate encoder learns entity-level
embeddings that condition the variable selection networks, gating, and
attention; the model thereby distinguishes entity-specific behaviour from
shared temporal structure. The original paper reports state-of-the-art
performance on a 369-store retail demand panel and the Volatility, Traffic,
Electricity, and Favorita benchmarks. To our knowledge TFT has not been
systematically applied to multi-market crisis EWS, an oversight likely
attributable to the assumption (implicit in much of the EWS literature) that
crisis prediction is best treated as a single-country regression. The static
embedding mechanism is precisely the architectural feature one would invoke
to share parameters across markets while preserving entity identity, and
exploiting it is the central reason TFT is the right base architecture for
this project.

\subsection{Regime Switching}
Regime-switching models have been a workhorse of financial econometrics
since \citet{hamilton1989new}, with extensive use in volatility modelling
\citep{angbekaert2002regime} and asset allocation. The classical
formulation is a low-dimensional Markov chain driving the parameters of an
otherwise linear model, fit by EM. In its modern, neural-network
incarnation, the regime structure becomes a latent variable in a deep
generative model; the forward algorithm remains differentiable and
end-to-end training becomes feasible. The principal limitation of classical
regime models for EWS is that they are typically estimated from a single
asset's returns and cannot exploit cross-market information. Our regime
module addresses both limitations: it operates on the (cross-market-aware)
LSTM encoder hidden state, and its transition matrix is itself a learned
function of the current state. The neural HMM literature in NLP and speech
\citep{tran2016unsupervised} establishes that this combination is trainable
in practice; the contribution here is its application to multi-market
financial crisis prediction in the regime-conditioned attention setup
described in Section~\ref{sec:method}.

\subsection{The Gap}
Three observations from the above set up the gap this project addresses.
Crisis EWS is a low-data regime when treated single-market, and tree models
dominate; but the data scarcity is a property of the problem formulation,
not of the underlying phenomenon, since multiple markets crash at different
times for related but distinct reasons. TFT is the natural architecture for
panel forecasting yet has not been applied to crisis EWS. Regime switching
is well established in finance but has not been integrated as a
differentiable layer inside a panel-forecasting attention model. RA-TFT sits
at the intersection of these three observations.

% =====================================================================
\section{Methodology}
\label{sec:method}

\subsection{Multi-Market Panel Reformulation}
The dataset comprises twelve equity indices: S\&P 500 ($\hat{}$GSPC),
S\&P/TSX ($\hat{}$GSPTSE), FTSE 100 ($\hat{}$FTSE), DAX ($\hat{}$GDAXI),
CAC 40 ($\hat{}$FCHI), Nikkei 225 ($\hat{}$N225), Hang Seng
($\hat{}$HSI), ASX 200 ($\hat{}$AXJO), KOSPI ($\hat{}$KS11), Bovespa
($\hat{}$BVSP), Nifty 50 ($\hat{}$NSEI), and Shanghai Composite
(000001.SS), each from 1 January 2000 to 31 December 2024. Indices are
synchronised to a global trading calendar with forward-fill of missing days
and dropped if more than 5\% of business days are missing in a calendar
year. The unit of observation is the (market, day) pair; static covariates
encode market identity. The panel contains roughly 30--50 distinct crisis
episodes (depending on the precise threshold) compared with the four
episodes typically isolated on the S\&P 500 alone, restoring a workable
data-to-parameter ratio for attention-based architectures.

\subsection{Labelling}
For each market $m$ and trading day $t$ we compute the trailing 252-day
rolling peak $P_{m,t}$ and the running drawdown
$D_{m,t} = P_{m,t}/X_{m,t} - 1$ where $X$ is the closing price. A market is
\emph{in crisis} on day $t$ if $D_{m,t} \geq 0.20$. The crisis onset day
is the first day of an in-crisis run preceded by at least twenty
non-crisis days. The binary target is
\(
y_{m,t} = \mathbf{1}\{\exists\,s \in [t+1, t+63]\text{ s.t. day $s$ is a crisis onset}\}.
\)
In-crisis observations are excluded from training to prevent the model from
learning trivial in-crisis dynamics.

\subsection{Features}
The historical input is a 252-day window terminating at $t$. Per-market
numeric features include log returns at 1/5/21-day horizons, realised
volatility at 5/21/63-day horizons, short-horizon drawdown from the 63-day
rolling peak, the 63-day Hurst exponent, vol-of-vol at 63 days, return
skewness at 63 days, a 63-day volume z-score, and Amihud illiquidity at 21
days. Global features (broadcast across all markets) include VIX, the
VIX-VIX3M term structure, US high-yield and investment-grade credit
spreads and their differential, the TED spread, the dollar index, gold, and
the MOVE index. Static categorical features comprise a single market
identity field with twelve levels, which is consumed by TFT's static
covariate encoder to produce a learned market embedding.

A previous iteration of the pipeline included a 252-day drawdown feature
that was found to encode the label definition almost perfectly; it was
removed and replaced with vol-of-vol and return skewness. The remaining
high-correlation features (VIX, credit spreads, dollar index) are genuine
economic precursors of crises, not leakage.

\subsection{RA-TFT Architecture}
RA-TFT extends the vanilla TFT by replacing two components. After the LSTM
encoder produces a hidden-state sequence
$\mathbf{h}_{1:T}\in\mathbb{R}^{T\times d}$, a \emph{regime detection
module} computes per-step regime posteriors
$\boldsymbol{\pi}_t \in \Delta^{K-1}$ over $K=3$ latent states using a
neural HMM:
\begin{equation}
\log p(z_t = k \mid \mathbf{h}_{1:t}) \;\propto\; \log p(\mathbf{h}_t \mid z_t=k) + \log\!\!\sum_{k'} A_{k'k}(\mathbf{h}_{t-1}) \, p(z_{t-1}=k' \mid \mathbf{h}_{1:t-1}),
\end{equation}
where the time-varying transition matrix $A(\mathbf{h})$ is a softmax
output of a small MLP. The transition MLP is initialised so that diagonal
entries dominate (sticky regimes), preventing oscillation in early
training. The emission distribution is implemented as a Gated Residual
Network mapping $\mathbf{h}_t$ to per-state log-likelihoods.

The inferred regime posteriors then condition an additive bias to the
attention logits in the interpretable multi-head attention block:
\begin{equation}
\text{logit}^{\text{(RA-TFT)}}_{ij} \;=\; \frac{\mathbf{q}_i^\top \mathbf{k}_j}{\sqrt{d}} + b_{\boldsymbol{\pi}_i, \boldsymbol{\pi}_j},
\end{equation}
where $b$ is a learned tensor of shape $K\times K\times H$ over (query
regime, key regime, attention head). This permits the attention pattern to
adapt to crisis regimes; for example, by attending more aggressively to
recent volatility spikes when the inferred state is `stress' than when it
is `calm'. The remaining TFT components (variable selection, static
encoders, gating, quantile head replaced by a sigmoid head for binary
classification) are inherited unchanged.

\subsection{Training}
The loss is binary cross-entropy with a class-frequency reweighting to
account for the $\sim$15\% positive rate. Optimisation uses AdamW
($\eta=10^{-3}$, weight decay $10^{-2}$), cosine annealing, and gradient
clipping at unit norm. Validation uses a random 20\% holdout of training
windows rather than a temporal cut, since temporal validation creates a
distribution shift between train and validation crisis rates that drives
spurious early stopping; this issue was diagnosed and resolved during
preliminary experimentation (see Section~\ref{sec:results}).

\subsection{Evaluation}
The principal metrics are ROC-AUC and PR-AUC; PR-AUC is the more
informative summary in this imbalanced regime. We report
\emph{pooled} scores (concatenating all markets and days in the test fold)
and \emph{macro-averaged} scores (mean of per-market scores), since pooling
weights observations by market activity while macro-averaging weights each
market equally. Calibration is assessed by expected calibration error and
Brier score, with post-hoc Platt and isotonic recalibration fit on the
validation set. Statistical significance uses bootstrap confidence
intervals on PR-AUC over the test fold.

The cross-validation protocol is crisis-anchored walk-forward: training
windows expand by one year per fold, and each test fold is required to
contain at least one crisis onset. This prevents the inflation of test
metrics that would occur on a quiet test period and ensures that the
reported numbers reflect the model's ability to anticipate crises rather
than identify their absence.

% =====================================================================
\section{Preliminary Results and Findings}
\label{sec:results}

The data pipeline, all four baseline models, and the RA-TFT architecture
have been implemented and exercised end-to-end on a single temporal split.
The principal results are summarised below; they are preliminary in the
sense that they use a single train/test cut rather than the full
walk-forward protocol planned for the dissertation, but they suffice to
establish feasibility and to validate the central hypothesis.

\subsection{Single-Market Baseline}
On the S\&P 500-only formulation with three crisis-anchored walk-forward
folds, the ranking of models is as one would predict from the prior
literature. Tree-based and linear models dominate sequence models:
\vspace{2pt}
\begin{center}
\small
\begin{tabular}{lcc}
\toprule
Model & Mean ROC-AUC & Mean PR-AUC \\
\midrule
XGBoost            & 0.655 & 0.507 \\
LSTM               & 0.529 & 0.309 \\
RA-TFT             & 0.484 & 0.310 \\
Vanilla TFT        & 0.477 & 0.281 \\
Logistic Regression & 0.403 & 0.294 \\
\bottomrule
\end{tabular}
\end{center}

XGBoost wins comfortably; the deep models are essentially indistinguishable
from random on ROC, consistent with \citet{bluwstein2023credit}. This
reproduces the literature consensus on single-market EWS.

\subsection{Multi-Market Panel}
On the twelve-market panel, with a fixed train/test split (training
2000--2018, testing 2019--2024) and the corrected validation protocol, the
ranking inverts:
\vspace{2pt}
\begin{center}
\small
\begin{tabular}{lcccc}
\toprule
Model & ROC-AUC & PR-AUC & ECE & Brier \\
\midrule
\textbf{RA-TFT}    & \textbf{0.640} & \textbf{0.220} & 0.329 & 0.265 \\
XGBoost            & 0.622 & 0.197 & 0.148 & 0.152 \\
Vanilla TFT        & 0.548 & 0.161 & 0.324 & 0.325 \\
LSTM               & 0.477 & 0.138 & 0.412 & 0.309 \\
\bottomrule
\end{tabular}
\end{center}

RA-TFT outperforms XGBoost on both ROC-AUC and PR-AUC, and outperforms
vanilla TFT by a substantial margin (0.640 vs.\ 0.548 ROC-AUC). This is
consistent with the central thesis: with the cross-entity pooling and the
regime module both engaged, the architectural advantages of attention-based
models for multi-entity panel forecasting are realised.

\subsection{Calibration}
RA-TFT's raw probabilities are over-confident (ECE $\approx 0.27$). Post-hoc
Platt scaling fitted on the validation set reduces this to 0.256 and Brier
score to 0.255. Isotonic regression performs slightly worse, consistent
with the limited validation set. XGBoost is already well-calibrated
($\text{ECE}=0.148$) and is barely affected by post-hoc methods. The
calibration gap between RA-TFT and XGBoost is the principal weakness of the
sequence model and a target for the remaining work.

\subsection{Lessons from Preliminary Work}
Three lessons have already shaped the methodology presented above. First,
the validation strategy has a first-order effect on attention-based
training: an early temporal-split protocol caused validation loss to
diverge from epoch one because the validation set contained almost no
crisis events, and early stopping consequently selected the least-trained
checkpoint. The current random-window holdout resolves this. Second, label
leakage through redundant drawdown features inflated single-market scores
by approximately ten ROC-AUC points before being identified and removed.
Third, sequence-model wins are extremely sensitive to seed and run-to-run
variance, with RA-TFT ROC-AUC observed in the range $0.570$--$0.640$
across seeds; multi-seed averaging is therefore mandatory in the final
walk-forward evaluation.

% =====================================================================
\section{Project Schedule}
\label{sec:schedule}

\subsection{Workstreams and Status}
The project decomposes into eight workstreams. Items already completed are
explicitly marked.

\begin{itemize}
\item[\textbf{W1.}] \emph{Topic selection \& literature review} (Dec 2025
-- Feb 2026). \textbf{Complete.}
\item[\textbf{W2.}] \emph{Data ingestion pipeline}: twelve indices, FRED
macro covariates, panel feature engineering, label generation
(Feb 2026). \textbf{Complete.}
\item[\textbf{W3.}] \emph{Baseline implementations}: XGBoost, logistic
regression, LSTM, vanilla TFT (Feb to Mar 2026). \textbf{Complete.}
\item[\textbf{W4.}] \emph{RA-TFT architecture}: regime module, regime
attention bias, end-to-end training (Mar to Apr 2026). \textbf{Complete}, with
single-split panel evaluation done.
\item[\textbf{W5.}] \emph{Walk-forward cross-validation on the panel}:
crisis-anchored expanding-window CV, multi-seed averaging, bootstrap CIs
(May to Jun 2026). Pending HPC allocation; cluster credits applied for
through King's e-Research.
\item[\textbf{W6.}] \emph{Ablation studies}: isolating the contribution
of the regime module, regime attention bias, and static market embeddings
(Jun 2026).
\item[\textbf{W7.}] \emph{Architectural extensions}: exploration of
either (a) a regime-routed mixture-of-experts on top of the existing regime
module, or (b) Mamba state-space backbone in place of the LSTM encoder. The
two options will be scoped in early June and one selected based on early
signal (Jun to Jul 2026).
\item[\textbf{W8.}] \emph{Per-market analysis, calibration, write-up,
video}: per-market ROC tables, reliability diagrams, dissertation
drafting, ten-minute video presentation (Jul to Aug 2026).
\end{itemize}

\subsection{Gantt Chart}
Figure~\ref{fig:gantt} maps the workstreams onto the King's MSc
Computational Finance project calendar. The four KCL milestones marked as
diamonds are: the present preliminary report submission (30 April 2026),
first individual feedback meeting (22 May 2026), second individual meeting
(17 July 2026), and final report and video submission (6 August 2026).
Group supervisions and progress slide deadlines impose intermediate
checkpoints which the workstream boundaries respect.

\begin{figure}[h]
\centering
\resizebox{\textwidth}{!}{%
\begin{ganttchart}[
  hgrid,
  vgrid,
  x unit=0.55cm,
  y unit chart=0.55cm,
  y unit title=0.55cm,
  bar/.append style={fill=black!55},
  bar incomplete/.append style={fill=black!25},
  milestone/.append style={fill=black},
  title label font=\scriptsize,
  bar label font=\scriptsize,
  milestone label font=\scriptsize,
  group label font=\scriptsize\bfseries,
]{1}{28}
  \gantttitle{2026 (calendar weeks from 1 Feb)}{28} \\
  \gantttitlelist{1,...,28}{1} \\
  \ganttbar{W1 Lit. review}{1}{4} \\
  \ganttbar{W2 Data pipeline}{1}{5} \\
  \ganttbar{W3 Baselines}{3}{8} \\
  \ganttbar{W4 RA-TFT build}{6}{13} \\
  \ganttmilestone{Prelim.\ report}{13} \\
  \ganttbar{W5 Walk-forward CV (HPC)}{14}{21} \\
  \ganttmilestone{1st individual feedback}{16} \\
  \ganttbar{W6 Ablations}{19}{22} \\
  \ganttbar{W7 Architectural extension}{20}{24} \\
  \ganttmilestone{2nd individual mtg}{24} \\
  \ganttbar{W8 Per-market \& calibration}{22}{26} \\
  \ganttbar{W8 Dissertation drafting}{23}{27} \\
  \ganttbar{W8 Video}{27}{27} \\
  \ganttmilestone{Final report \& video}{27} \\
\end{ganttchart}}
\caption{Project Gantt chart, weeks measured from the start of work on
1 February 2026. Diamond markers denote King's MSc Computational Finance
milestones; bars denote the eight project workstreams. Workstreams W1--W4
are complete at the time of writing.}
\label{fig:gantt}
\end{figure}

\subsection{Risks and Mitigations}
The principal external dependency is compute. RunPod credits (used for the
preliminary panel experiments) have proved expensive at scale, and the
walk-forward CV on the panel is estimated to require $\sim$200 GPU-hours
across folds and seeds. An application has been submitted for an HPC
allocation through King's e-Research, which would absorb essentially all of
W5 and W7. Should this allocation be delayed, the contingency plan is to
reduce the seed count from five to three and the walk-forward fold count
from eight to five, which preserves the statistical protocol at the cost of
slightly wider confidence intervals.

A secondary risk is that the run-to-run variance observed in preliminary
RA-TFT runs (ROC-AUC $0.57$--$0.64$ across seeds) does not tighten under
multi-seed averaging, in which case the headline claim of RA-TFT $>$
XGBoost would not be statistically significant despite a positive mean
effect. The mitigation here is structural: the ablation in W6 will isolate
which component(s) of RA-TFT contribute robustly, allowing a more
defensible secondary claim even if the headline comparison is not
significant. A tertiary risk is overfitting to the present 2000--2024
window; a small held-out 2025 stress in any of the twelve markets (such as
intermittent Chinese property-sector volatility) would provide an external
validity check, and the pipeline already supports it.

% =====================================================================
\section{Conclusion}
\label{sec:conclusion}

This project investigates a specific and previously underexploited
hypothesis: that the apparent inferiority of attention-based sequence
models in financial crisis early warning is a consequence of single-market
data scarcity rather than a fundamental property of the architecture. The
RA-TFT proposal (a Temporal Fusion Transformer augmented with a neural
HMM regime detection module and regime-conditioned attention bias, applied
to a twelve-market equity panel) is designed to test this hypothesis at
the smallest scale where the test is informative. Preliminary
single-split results confirm the central prediction: on the panel
formulation, RA-TFT $0.640 >$ XGBoost $0.622$ in pooled ROC-AUC, reversing
the literature consensus on single-market EWS. The remaining work between
the present preliminary report and the 6 August 2026 final submission
concerns three things: replacing the single split with a crisis-anchored
walk-forward protocol, ablating the architectural components of RA-TFT to
identify which drives the improvement, and one architectural extension
(mixture-of-experts or Mamba backbone) to deepen the methodological
contribution. The schedule for these is set out above.

% =====================================================================
% References (do not count toward the 10-page body limit)
% =====================================================================
\begin{thebibliography}{99}

\bibitem[Ang and Bekaert(2002)]{angbekaert2002regime}
Ang, A.\ and Bekaert, G.\ (2002).
\newblock Regime switches in interest rates.
\newblock \emph{Journal of Business and Economic Statistics}, 20(2):163--182.

\bibitem[Beutel et al.(2019)]{beutel2019does}
Beutel, J., List, S.\ and von Schweinitz, G.\ (2019).
\newblock Does machine learning help us predict banking crises?
\newblock \emph{Journal of Financial Stability}, 45:100693.

\bibitem[Bluwstein et al.(2023)]{bluwstein2023credit}
Bluwstein, K., Buckmann, M., Joseph, A., Kapadia, S.\ and \c{S}im\c{s}ek, \"O.\ (2023).
\newblock Credit growth, the yield curve and financial crisis prediction:
evidence from a machine learning approach.
\newblock \emph{Journal of International Economics}, 145:103773.

\bibitem[Chatzis et al.(2018)]{chatzis2018forecasting}
Chatzis, S.\ P., Siakoulis, V., Petropoulos, A., Stavroulakis, E.\ and Vlachogiannakis, N.\ (2018).
\newblock Forecasting stock market crisis events using deep and statistical
machine learning techniques.
\newblock \emph{Expert Systems with Applications}, 112:353--371.

\bibitem[Dichtl et al.(2023)]{dichtl2023ews}
Dichtl, H., Drobetz, W.\ and Otto, T.\ (2023).
\newblock Forecasting stock market crashes via machine learning.
\newblock \emph{Journal of Financial Stability}, 65:101099.

\bibitem[Frankel and Saravelos(2012)]{frankelsaravelos2012}
Frankel, J.\ and Saravelos, G.\ (2012).
\newblock Can leading indicators assess country vulnerability? Evidence from
the 2008--09 global financial crisis.
\newblock \emph{Journal of International Economics}, 87(2):216--231.

\bibitem[Hamilton(1989)]{hamilton1989new}
Hamilton, J.\ D.\ (1989).
\newblock A new approach to the economic analysis of nonstationary time
series and the business cycle.
\newblock \emph{Econometrica}, 57(2):357--384.

\bibitem[Holopainen and Sarlin(2017)]{holopainen2017toward}
Holopainen, M.\ and Sarlin, P.\ (2017).
\newblock Toward robust early-warning models: a horse race, ensembles and
model uncertainty.
\newblock \emph{Quantitative Finance}, 17(12):1933--1963.

\bibitem[Kaminsky and Reinhart(1999)]{kaminsky1999leading}
Kaminsky, G.\ L.\ and Reinhart, C.\ M.\ (1999).
\newblock The twin crises: the causes of banking and balance-of-payments
problems.
\newblock \emph{American Economic Review}, 89(3):473--500.

\bibitem[Lim et al.(2021)]{lim2021tft}
Lim, B., Ar{\i}k, S.\ \"O., Loeff, N.\ and Pfister, T.\ (2021).
\newblock Temporal Fusion Transformers for interpretable multi-horizon time
series forecasting.
\newblock \emph{International Journal of Forecasting}, 37(4):1748--1764.

\bibitem[Reinhart and Rogoff(2009)]{reinhartrogoff2009}
Reinhart, C.\ M.\ and Rogoff, K.\ S.\ (2009).
\newblock \emph{This Time Is Different: Eight Centuries of Financial Folly}.
\newblock Princeton University Press.

\bibitem[Tran et al.(2016)]{tran2016unsupervised}
Tran, K., Bisk, Y., Vaswani, A., Marcu, D.\ and Knight, K.\ (2016).
\newblock Unsupervised neural hidden Markov models.
\newblock In \emph{Proceedings of the Workshop on Structured Prediction for
NLP}, EMNLP, pages 63--71.

\end{thebibliography}

\end{document}
